\section{The torus representation and the local law}\label{app:torus}

This appendix proves Lemma~\ref{lem:upper}, Theorem~\ref{thm:local} and
Lemma~\ref{lem:discrete}. All 3 rest on a formula that writes the run sum
$H(t)$ of Lemma~\ref{lem:continuum} as an integral over a torus.

Throughout, $F$ is a face with $k\ge2$ states, numbered $1,\dots,k$. For
$x\in\mathbb C^F$ we write $x'=(x_2,\dots,x_k)$, and $\ones$ is the all-ones
vector of the appropriate size. We write $\mathrm i$ for the imaginary unit,
since $i$ denotes a state, and $\mathbb T^m=(\R/2\pi\Z)^m$. We split
\[
A_F=\begin{pmatrix}0&a_1^\top\\ b_1&A'\end{pmatrix},\qquad
\mu_F=(\mu_1,\mu'),
\]
so $a_1=(q_{1j})_{j\ge2}$, $b_1=(q_{i1})_{i\ge2}$ and
$A'=(q_{ij})_{i,j\ge2}$. For $z'\in\mathbb C^{k-1}$ let
$Z'=\operatorname{diag}(z')$. Where $Z'-A'$ is invertible we put
\[
X(z')=(Z'-A')^{-1},\qquad g(z')=a_1^\top X(z')b_1,\qquad
b(z')=\bigl(\mu_1+\mu'^\top X(z')b_1\bigr)\bigl(1+a_1^\top X(z')\ones\bigr).
\]
We use 2 standard facts. First, let $B\ge0$ be a square matrix and $S$ a
diagonal matrix with positive diagonal. Then $\varrho(B-S)<0$ holds if and
only if $\varrho(S^{-1}B)<1$, and also if and only if $(S-B)v>0$ for some
vector $v>0$. In that case
$(S-B)^{-1}=\sum_{j\ge0}(S^{-1}B)^jS^{-1}\ge0$. These are among the standard
characterizations of nonsingular M-matrices~\cite[Ch.~6]{bermanplemmons1994}.
Second, if $M$ is an irreducible Metzler matrix, then $\varrho(M)$ is larger
than the Perron root of each proper principal submatrix of $M$, and every
nonnegative eigenvector of $M$ is positive and belongs to
$\varrho(M)$~\cite{bermanplemmons1994}.

\begin{lemma}[torus representation]\label{lem:torus}
Let $s'\in(0,\infty)^{k-1}$ satisfy
$\varrho(A'-\operatorname{diag}s')<0$. On the set of $z'\in\mathbb C^{k-1}$
with $\lvert z_i\rvert\ge s_i$ for every $i\ge2$, the matrix $Z'-A'$ is
invertible, and $X$, $g$ and $b$ are given by absolutely convergent power
series in $1/z_2,\dots,1/z_k$ with nonnegative coefficients. In particular
$\lvert g(z')\rvert\le g(s')$ and $\lvert b(z')\rvert\le b(s')$ there. For
every $t\in(0,\infty)^F$,
\[
H(t)=\frac1{(2\pi)^{k-1}}\int_{\mathbb T^{k-1}}
\exp\Bigl(t_1g(z')+\sum_{i\ge2}t_iz_i\Bigr)\,b(z')\prod_{i\ge2}z_i\,d\varphi,
\qquad z_i=s_ie^{\mathrm i\varphi_i}.
\]
\end{lemma}

\begin{proof}
Put $S'=\operatorname{diag}(s')$. By the first fact,
$\varrho(S'^{-1}A')<1$. So the series
\[
X(z')=\sum_{j\ge0}(Z'^{-1}A')^jZ'^{-1}
\]
converges absolutely for $\lvert z_i\rvert\ge s_i$, since its terms are
bounded entrywise in modulus by those of the same series at $z'=s'$. This
gives the claims about $X$, $g$ and $b$, because $a_1$, $b_1$ and $\mu$ are
nonnegative.

Choose $s_1>g(s')$, put $s=(s_1,s')$ and $S=\operatorname{diag}(s)$. We show
that $\varrho(S^{-1}A_F)<1$. This is the Schur complement criterion for
M-matrices, and we check it directly. For $\varepsilon>0$ let $v=(1,v')$
with $v'=X(s')(b_1+\varepsilon\ones)$. Then $v'\ge
S'^{-1}(b_1+\varepsilon\ones)>0$, and
\[
\bigl((S-A_F)v\bigr)_i=\bigl((S'-A')v'\bigr)_i-(b_1)_i=\varepsilon\quad(i\ge2),
\qquad
\bigl((S-A_F)v\bigr)_1=s_1-g(s')-\varepsilon\,a_1^\top X(s')\ones .
\]
For small $\varepsilon$ both are positive, so the first fact gives the
claim.

Now let $\mathcal Z=\operatorname{diag}(z)$ with $\lvert z_i\rvert=s_i$ for
every $i\in F$. Expanding
$(\mathcal Z-A_F)^{-1}=\mathcal Z^{-1}(I-A_F\mathcal Z^{-1})^{-1}$ in powers
of $A_F\mathcal Z^{-1}$ gives
\[
\mu_F^\top(\mathcal Z-A_F)^{-1}\ones=\sum_c\mu_{c_1}\prod_{j<M}q_{c_jc_{j+1}}
\prod_{i\in F}z_i^{-m_i(c)},
\]
where $c=c_1\cdots c_M$ runs over all words in $F$ with $M\ge1$ and
$c_j\ne c_{j+1}$, and $m_i(c)$ is the number of letters $i$ in $c$. Since
$\varrho(A_FS^{-1})=\varrho(S^{-1}A_F)<1$, the expansion is valid and the
sum converges absolutely and uniformly on the torus. For an integer $m$, $t_i>0$ and
$z_i=s_ie^{\mathrm i\varphi_i}$,
\[
\frac1{2\pi}\int_0^{2\pi}e^{t_iz_i}z_i^{1-m}\,d\varphi_i
=\frac1{2\pi\mathrm i}\oint_{\lvert z_i\rvert=s_i}e^{t_iz_i}z_i^{-m}\,dz_i
=\begin{cases}t_i^{m-1}/(m-1)!,&m\ge1,\\ 0,&m\le0.\end{cases}
\]
We integrate term by term. The words that miss a state of $F$ give $0$, and
grouping the others by $m=(m_i(c))_{i\in F}$ gives
$W(m)\prod_it_i^{m_i-1}/(m_i-1)!$. Hence
\begin{equation}\label{eq:A-full}
H(t)=\frac1{(2\pi)^k}\int_{\mathbb T^k}e^{\langle t,z\rangle}\,
\mu_F^\top(\mathcal Z-A_F)^{-1}\ones\prod_{i\in F}z_i\,d\varphi .
\end{equation}
We do the integral over $\varphi_1$ first. Put $\zeta=z_1-g(z')$, the Schur
complement of $Z'-A'$ in $\mathcal Z-A_F$. It is not $0$, because
$\lvert g(z')\rvert\le g(s')<s_1$. With $X=X(z')$, block inversion gives
$(\mathcal Z-A_F)^{-1}$ with first row $\zeta^{-1}(1,a_1^\top X)$, first
column $\zeta^{-1}(1,Xb_1)^\top$ and lower right block
$X+\zeta^{-1}Xb_1a_1^\top X$. So
\[
\mu_F^\top(\mathcal Z-A_F)^{-1}\ones=\frac{b(z')}{z_1-g(z')}+\mu'^\top X(z')\ones .
\]
The second term does not depend on $z_1$. So by Cauchy's integral formula the integral over
$\varphi_1$ in~\eqref{eq:A-full}, divided by $2\pi$, equals
$e^{t_1g(z')}b(z')$ times the factors that do not involve $z_1$. In other
words $b(z')$ is the residue at the pole $z_1=g(z')$. This proves the
formula.
\end{proof}

The reason for this form is positivity. All coefficients are nonnegative, so
on the torus the modulus of the integrand is at most
$\exp(t_1g(s')+\sum_{i\ge2}t_is_i)\,b(s')\prod_{i\ge2}s_i$, which is its
value at $\varphi=0$. This bound gives Lemma~\ref{lem:upper}, and a Laplace
expansion around $\varphi=0$ gives Theorem~\ref{thm:local}.

\subsection{Proof of Lemma~\texorpdfstring{\ref{lem:upper}}{3.2}}

Let $F$ be irreducible and $\theta\in V_F$. Then $A_F-\operatorname{diag}\theta$
is an irreducible Metzler matrix with Perron root $0$. So it has a right
eigenvector $u>0$ for the eigenvalue $0$, and
$\theta_i=(A_Fu)_i/u_i$. Since $G_F$ is strongly connected and $k\ge2$, each
row of $A_F$ has a positive entry, so $\theta>0$. By the second fact,
$\varrho(A'-\operatorname{diag}\theta')<0$. The rows $i\ge2$ of
$(A_F-\operatorname{diag}\theta)u=0$ read
$(\operatorname{diag}\theta'-A')u'=u_1b_1$, so $u'=u_1X(\theta')b_1$. Then
row $1$ gives $\theta_1u_1=a_1^\top u'=u_1g(\theta')$. Hence
\begin{equation}\label{eq:A-radii}
\theta>0,\qquad \varrho(A'-\operatorname{diag}\theta')<0,\qquad
g(\theta')=\theta_1\qquad\text{for every }\theta\in V_F .
\end{equation}

Now fix $n$ with support $F$ and $h$ with $hq_{\max}\le\frac12$, and put
$\lvert n\rvert=\sum_in_i$ and $t=hn$. Lemma~\ref{lem:runs}, applied with
$\lvert n\rvert$ in place of $N$ and the same $h$, gives the probability
$P_{\lvert n\rvert}(n)$ that the first $\lvert n\rvert$ visits have
composition $n$. Only terms with $m_i\le n_i$ for every $i$ are nonzero. In
them we use $h^{m_i-1}\binom{n_i-1}{m_i-1}\le t_i^{m_i-1}/(m_i-1)!$ and
$(1-hq_i)^{n_i-m_i}\le e^{-q_it_i}\sigma^{m_i}$, where
$\sigma=e^{hq_{\max}}$. Note that $h^{M-1}=h^{k-1}\prod_ih^{m_i-1}$, and
that $\sigma^MW(m)=\sigma W_\sigma(m)$, where $W_\sigma$ is $W$
computed with $A_F$ replaced by $\sigma A_F$. So
\begin{equation}\label{eq:A-sigma}
P_{\lvert n\rvert}(n)\le\sigma\,h^{k-1}e^{-\langle q,t\rangle}H_\sigma(t),
\end{equation}
where $H_\sigma$ is $H$ computed with $\sigma A_F$ in place of $A_F$.
We bound $H_\sigma(t)$ by Lemma~\ref{lem:torus} for $\sigma A_F$ with the
radii $\sigma\theta'$. They satisfy its hypothesis, since
$\varrho(\sigma A'-\operatorname{diag}\sigma\theta')
=\sigma\varrho(A'-\operatorname{diag}\theta')<0$. The functions of
Lemma~\ref{lem:torus} for $\sigma A_F$ satisfy
$X_\sigma(\sigma z')=\sigma^{-1}X(z')$, so
$g_\sigma(\sigma z')=\sigma g(z')$ and
$b_\sigma(\sigma z')=b(z')$. By~\eqref{eq:A-radii} and the bounds in
Lemma~\ref{lem:torus}, the integrand for $\sigma A_F$ has modulus at most
$e^{\sigma\langle t,\theta\rangle}b(\theta')\sigma^{k-1}\prod_{i\ge2}\theta_i$.
With~\eqref{eq:A-sigma} this gives
\[
P_{\lvert n\rvert}(n)\le\sigma^kb(\theta')\prod_{i\ge2}\theta_i\;
h^{k-1}\exp\bigl(-h\langle n,q-\theta\rangle+(\sigma-1)h\langle n,\theta\rangle\bigr).
\]
Finally $\sigma\le e^{1/2}$ and $\sigma-1\le2hq_{\max}$, so
$(\sigma-1)h\langle n,\theta\rangle\le2q_{\max}(\max_i\theta_i)h^2\lvert n\rvert$.
This proves the lemma with
$C_\theta=\max\bigl(e^{k/2}b(\theta')\prod_{i\ge2}\theta_i,\ 2q_{\max}\max_i\theta_i\bigr)$.
For the bound stated after the lemma take
$\theta=q-\lambda_F\ones$. It lies in $V_F$ because
$\varrho(A_F-\operatorname{diag}\theta)=\varrho(Q_F)+\lambda_F=0$. Then
$\langle n,q-\theta\rangle=\lambda_F\lvert n\rvert$, and for
$\lvert n\rvert=N$ we have $h\lvert n\rvert=T$ and
$h^2\lvert n\rvert=T^2/N$. \qed

\subsection{Proof of Theorem~\texorpdfstring{\ref{thm:local}}{4.3}}

For $k=1$ the formula is part of Lemma~\ref{lem:continuum}. Let $k\ge2$, fix
$\alpha\in\Delta_F^\circ$, and write $\theta=\theta(\alpha)$, $r=r_\alpha$
and $l=l_\alpha$ as in Definition~\ref{def:tilt}, and
$\Theta'=\operatorname{diag}(\theta')$. By~\eqref{eq:A-radii} we may use
Lemma~\ref{lem:torus} with the radii $\theta'$ and $t=T\alpha$. With
Lemma~\ref{lem:continuum} this gives
\begin{equation}\label{eq:A-fT}
f_T(\alpha)=\frac{T^{k-1}e^{-T\langle q,\alpha\rangle}}{(2\pi)^{k-1}}
\int_{\mathbb T^{k-1}}e^{T\Psi(\alpha,\varphi)}\chi(\alpha,\varphi)\,d\varphi,
\end{equation}
where $z_i=\theta_ie^{\mathrm i\varphi_i}$ and
\[
\Psi(\alpha,\varphi)=\alpha_1g(z')+\sum_{i\ge2}\alpha_iz_i,\qquad
\chi(\alpha,\varphi)=b(z')\prod_{i\ge2}z_i .
\]
We check the hypotheses of Lemma~\ref{lem:laplace} with $m=k-1$, phase $\Psi$
and amplitude $\chi$, and then compute the constant.

We first show that $\varphi=0$ is the only maximum of
$\operatorname{Re}\Psi(\alpha,\cdot)$ on $\mathbb T^{k-1}$.
By~\eqref{eq:A-radii}, $\Psi(\alpha,0)=\alpha_1\theta_1+\sum_{i\ge2}\alpha_i\theta_i
=\langle\alpha,\theta\rangle$, which is real. By Lemma~\ref{lem:torus},
$\lvert g(z')\rvert\le g(\theta')=\theta_1$ on the torus, so
\[
\operatorname{Re}\Psi(\alpha,\varphi)\le\alpha_1\theta_1+\sum_{i\ge2}\alpha_i\theta_i\cos\varphi_i
=\Psi(\alpha,0)-\sum_{i\ge2}\alpha_i\theta_i(1-\cos\varphi_i).
\]
Every $\alpha_i$ is positive, and every $\theta_i$ is positive
by~\eqref{eq:A-radii}. So the last sum is positive unless $\varphi=0$ in
$\mathbb T^{k-1}$, which is hypothesis (iii) of Lemma~\ref{lem:laplace}. Note
that the gain comes from the angles and not from $g$. For $k=2$,
$g(z_2)=q_{12}q_{21}/z_2$ has constant modulus on the circle. No
aperiodicity condition is needed.

Next we compute the derivatives at $\varphi=0$. For $z\in\R^F$ near $\theta$
let $\lambda(z)=\varrho(A_F-\operatorname{diag}z)$, and for $\eta\in\R^F$ near
$0$ let $\Lambda(\eta)=\varrho(\widehat Q_\alpha+\operatorname{diag}\eta)$.
Since $A_F-\operatorname{diag}(\theta-\eta)
=R(\widehat Q_\alpha+\operatorname{diag}\eta)R^{-1}$ with
$R=\operatorname{diag}r$, we have $\lambda(\theta-\eta)=\Lambda(\eta)$. The
Perron root of the irreducible matrix $\widehat Q_\alpha+\operatorname{diag}\eta$
is a simple eigenvalue. So by the implicit function theorem, applied to the
characteristic polynomial, $\Lambda$ and the right eigenvector $v(\eta)$ with
$\alpha^\top v(\eta)=1$ are real-analytic near $0$, with $\Lambda(0)=0$ and
$v(0)=\ones$. Recall $Z\ones=0$, $\alpha^\top Z=0$ and
$-\widehat Q_\alpha Z=I-\ones\alpha^\top$ for the matrix $Z$ of
Definition~\ref{def:tilt}. These follow from
$(\ones\alpha^\top-\widehat Q_\alpha)\ones=\ones$ and
$\alpha^\top(\ones\alpha^\top-\widehat Q_\alpha)=\alpha^\top$. Differentiating
$(\widehat Q_\alpha+\operatorname{diag}\eta)v=\Lambda v$ at $\eta=0$ gives
$e_j+\widehat Q_\alpha\partial_jv=(\partial_j\Lambda)\ones$. Multiplying by
$\alpha^\top$ gives $\partial_j\Lambda=\alpha_j$. Then $\partial_jv=Ze_j$,
because both vectors solve $\widehat Q_\alpha x=\alpha_j\ones-e_j$ with
$\alpha^\top x=0$, and the kernel of $\widehat Q_\alpha$ is spanned by $\ones$.
Differentiating twice and multiplying by $\alpha^\top$ gives
$\partial_i\partial_j\Lambda(0)=\alpha_iZ_{ij}+\alpha_jZ_{ji}$. Hence
\[
\nabla\lambda(\theta)=-\alpha,\qquad\nabla^2\lambda(\theta)=\Gamma_\alpha .
\]
For real $z'$ near $\theta'$ we have $\varrho(A'-\operatorname{diag}z')<0$ by
continuity, so the vector $(1,X(z')b_1)$ is nonnegative by the first fact. The
definition of $g$ shows that this vector lies in the kernel of
$A_F-\operatorname{diag}(g(z'),z')$, so $\lambda(g(z'),z')=0$ by the second
fact. Write $g_j$ and $g_{ij}$ for the first and second
partial derivatives of $g$ at $\theta'$, and put
$w_j=e_j-(\alpha_j/\alpha_1)e_1\in\R^F$ for $j\ge2$. Differentiating
$\lambda(g(z'),z')=0$ once and twice at $z'=\theta'$ gives
\[
g_j=-\frac{\alpha_j}{\alpha_1},\qquad
\alpha_1g_{ij}=w_i^\top\Gamma_\alpha w_j .
\]
So $\partial\Psi/\partial z_j=\alpha_1g_j+\alpha_j=0$ at $\theta'$, and the
matrix $\mathcal H=(\partial^2\Psi/\partial z_i\partial z_j)_{i,j\ge2}$ at
$\theta'$ equals $V^\top\Gamma_\alpha V$, where $V$ has columns
$w_2,\dots,w_k$. Since $\partial/\partial\varphi_j=\mathrm iz_j\,\partial/\partial z_j$,
at $\varphi=0$ we get $\nabla_\varphi\Psi=0$ and
$\nabla_\varphi^2\Psi=-\Theta'\mathcal H\Theta'$, because the extra term
involves only first derivatives, which vanish there. So hypothesis (i) of
Lemma~\ref{lem:laplace} holds, and $-\nabla^2_\varphi\Psi(\alpha,0)=\Theta'\mathcal H\Theta'$
is real and symmetric.

Now we show that $\mathcal H$ is positive definite and find its determinant.
Let $\Sigma_\alpha$ be $\Gamma_\alpha$ without the row and column of state
$1$, and $B=I+\ones\alpha'^\top/\alpha_1$. For $y\in\R^{k-1}$ we have
$Vy+(\alpha'^\top y/\alpha_1)\ones=(0,By)$. Since $\Gamma_\alpha\ones=0$,
adding a multiple of $\ones$ does not change the quadratic form, so
$y^\top\mathcal Hy=(By)^\top\Sigma_\alpha(By)$. The matrix determinant lemma
gives $\det B=1+(1-\alpha_1)/\alpha_1=1/\alpha_1$. Hence
\[
\det\mathcal H=\frac{\det\Sigma_\alpha}{\alpha_1^2},
\]
and $\mathcal H$ is positive definite if $\Sigma_\alpha$ is. Let
$y\in\R^{k-1}$ be nonzero and $x=(0,y)\in\R^F$. Put
$f=x-(\alpha^\top x)\ones$ and $u=Zf$. Then $\alpha^\top f=0$,
$\alpha^\top u=0$ and $-\widehat Q_\alpha u=f$. With
$D=\operatorname{diag}\alpha$, using $\Gamma_\alpha\ones=0$ again,
\[
y^\top\Sigma_\alpha y=f^\top\Gamma_\alpha f=2f^\top DZf
=2\sum_i\alpha_iu_i(-\widehat Q_\alpha u)_i
=\sum_{i\ne j}\alpha_i(\widehat Q_\alpha)_{ij}(u_i-u_j)^2 .
\]
The last step uses $\widehat Q_\alpha\ones=0$ and
$\alpha^\top\widehat Q_\alpha=0$. For $i\ne j$,
$(\widehat Q_\alpha)_{ij}=q_{ij}r_j/r_i$ is positive exactly on the edges of
$G_F$, which is strongly connected. So the sum vanishes only if $u$ is
constant. Then $u=0$ since $\alpha^\top u=0$, so $f=0$ and $x$ is a multiple
of $\ones$. This is impossible, since $x_1=0$ and $x\ne0$. Hence
$\Sigma_\alpha$ and $\mathcal H$ are positive definite, and hypothesis (ii)
of Lemma~\ref{lem:laplace} holds.

Next we check the regularity. First, $\alpha\mapsto\theta(\alpha)$ is
real-analytic on $\Delta_F^\circ$. Indeed, since
$I_F(\alpha)=\langle q,\alpha\rangle-\inf_{u>0}\sum_i\alpha_i(A_Fu)_i/u_i$,
Proposition~\ref{prop:rate}(b) says that $r_\alpha$ is the unique minimizer
of this function up to scaling. With $u=e^w$ and $w_1=0$ the function is
$\mathcal U(\alpha,w)=\sum_{i\ne j}\alpha_iq_{ij}e^{w_j-w_i}$. Its Hessian in
$(w_2,\dots,w_k)$ is the quadratic form
$y\mapsto\sum_{i\ne j}\alpha_iq_{ij}e^{w_j-w_i}(y_j-y_i)^2$ with $y_1=0$,
which is positive definite because $G_F$ is strongly connected. So the
analytic implicit function theorem, applied to
$\nabla_w\mathcal U(\alpha,w)=0$, shows that $r_\alpha$ and
$\theta(\alpha)=A_Fr_\alpha/r_\alpha$ are real-analytic in $\alpha$. Now fix
$\alpha_0\in\mathcal K$. By~\eqref{eq:A-radii} and continuity of the Perron
root there is $\delta\in(0,\min_i\theta_i(\alpha_0))$ with
$\varrho(A'-\operatorname{diag}(\theta'(\alpha_0)-\delta\ones))<0$. By
Lemma~\ref{lem:torus}, $X$, $g$ and $b$ are holomorphic on the set where
$\lvert z_i\rvert>\theta_i(\alpha_0)-\delta$ for every $i\ge2$, as locally
uniform limits of their series. For $\alpha$ near $\alpha_0$ we have
$\theta_i(\alpha)>\theta_i(\alpha_0)-\delta$. So $\Psi$ and $\chi$ are
real-analytic in $(\alpha,\varphi)$ on a neighborhood of
$\{\alpha_0\}\times\mathbb T^{k-1}$. Covering $\mathcal K$ by finitely many
such neighborhoods, we see that $\Psi$ and $\chi$ are jointly real-analytic
near $\mathcal K\times\mathbb T^{k-1}$, and all their $\varphi$-derivatives
are jointly continuous on $\mathcal K\times\mathbb T^{k-1}$.

Lemma~\ref{lem:laplace} now gives, uniformly for $\alpha\in\mathcal K$ and
$T\ge1$,
\[
\int_{\mathbb T^{k-1}}e^{T\Psi}\chi\,d\varphi
=e^{T\langle\alpha,\theta\rangle}\Bigl(\frac{2\pi}T\Bigr)^{\frac{k-1}2}
\det(\Theta'\mathcal H\Theta')^{-1/2}
\Bigl[\chi(\alpha,0)+\frac{c_1(\alpha)}T+O(T^{-3/2})\Bigr],
\]
where $c_1$ is the coefficient given there. Here
$\det(\Theta'\mathcal H\Theta')=(\prod_{i\ge2}\theta_i)^2\det\Sigma_\alpha/\alpha_1^2$
and $\chi(\alpha,0)=b(\theta')\prod_{i\ge2}\theta_i$. To evaluate
$b(\theta')$, which is the residue of the generating function at
$z_1=g(\theta')$, we use the Perron vectors. As in the derivation
of~\eqref{eq:A-radii}, $r'=r_1X(\theta')b_1$. In the same way the columns
$j\ge2$ of $l^\top(A_F-\operatorname{diag}\theta)=0$ give
$l'=l_1X(\theta')^\top a_1$. Hence $\mu\cdot r=r_1(\mu_1+\mu'^\top X(\theta')b_1)$
and $l\cdot\ones=l_1(1+a_1^\top X(\theta')\ones)$. Since $l_1r_1=\alpha_1$,
\[
b(\theta')=\frac{(\mu\cdot r_\alpha)(l_\alpha\cdot\ones)}{\alpha_1},
\qquad\text{so}\qquad
\frac{\chi(\alpha,0)}{\det(\Theta'\mathcal H\Theta')^{1/2}}
=\frac{(\mu\cdot r_\alpha)(l_\alpha\cdot\ones)}{(\det\Sigma_\alpha)^{1/2}} .
\]
Here we use $\mu(F)>0$. Since $r_\alpha,l_\alpha>0$, the value
$\chi(\alpha,0)$ is positive exactly when $\mu(F)>0$, which holds by
\textup{(H)}. It is continuous in $\alpha$, so it is bounded below on
$\mathcal K$, and we may divide by it. (If $\mu(F)=0$, then $f_T=0$ and the
relative form makes no sense.) With $c(\alpha)=c_1(\alpha)/\chi(\alpha,0)$,
equation~\eqref{eq:A-fT} becomes
\[
f_T(\alpha)=T^{\frac{k-1}2}(2\pi)^{-\frac{k-1}2}
\frac{(\mu\cdot r_\alpha)(l_\alpha\cdot\ones)}{(\det\Sigma_\alpha)^{1/2}}\,
e^{-T\langle q-\theta,\alpha\rangle}\Bigl(1+\frac{c(\alpha)}T+O(T^{-3/2})\Bigr),
\]
uniformly on $\mathcal K$. Since $I_F(\alpha)=\langle q-\theta(\alpha),\alpha\rangle$,
this is the claim. The function $c$ is real-analytic, since $c_1$ is by
Lemma~\ref{lem:laplace} and $\chi(\alpha,0)>0$ is. It is real, because
$f_T(\alpha)$ and the main term are real, so
$\operatorname{Im}c(\alpha)=O(T^{-1/2})$ for every $T\ge1$, which forces
$\operatorname{Im}c(\alpha)=0$. \qed

\subsection{Proof of Lemma~\texorpdfstring{\ref{lem:discrete}}{4.4}}

We compare Lemma~\ref{lem:runs} with its continuum version term by term. Under the continuum weights a path has at most of order $T$ runs,
and each run after the first in state $i$ changes the number of placements by
a factor $1-O(m_i/n_i)$. So the relative error is of order $T^2/N$.

For $k=1$ we have $P_N(Ne_i)=\mu_i(1-hq_i)^{N-1}$ and $f_T=\mu_ie^{-q_iT}$.
Since $hq_i\le\frac12$, $(N-1)\log(1-hq_i)+q_iT=O(T/N+T^2/N)=O(T^2/N)$, and
$T^2/N\le1$, so the ratio is $1+O(T^2/N)$. Let
$k\ge2$, and write $t=hn$ and $\alpha=n/N\in\mathcal K$. By
Lemma~\ref{lem:continuum},
$N^{-(k-1)}f_T(n/N)=h^{k-1}e^{-\langle q,t\rangle}H(t)=\sum_mv_m$ with
\[
v_m=W(m)\,h^{k-1}\prod_{i\in F}e^{-q_it_i}\frac{t_i^{m_i-1}}{(m_i-1)!},
\qquad m\in\Z_{\ge1}^F .
\]
Since $h^{m_i-1}\binom{n_i-1}{m_i-1}=\frac{t_i^{m_i-1}}{(m_i-1)!}\prod_{1\le j<m_i}(1-j/n_i)$
and $h^{M-1}=h^{k-1}\prod_ih^{m_i-1}$, the $m$-th term of
Lemma~\ref{lem:runs} is $\kappa_mv_m$ with
\[
\kappa_m=\prod_{i\in F}\Bigl[\prod_{1\le j<m_i}\Bigl(1-\frac j{n_i}\Bigr)\Bigr]
(1-hq_i)^{n_i-m_i}e^{hq_in_i}.
\]
We bound $\kappa_m$. Since $(1-hq_i)^{n_i-m_i}\le e^{-hq_i(n_i-m_i)}$, we
get $\kappa_m\le e^{hq_{\max}M}$. For the lower bound,
$\log(1-x)\ge-x-x^2$ for $0\le x\le\frac12$ gives
$(1-hq_i)^{n_i-m_i}e^{hq_in_i}\ge(1-hq_i)^{n_i}e^{hq_in_i}\ge e^{-n_ih^2q_i^2}$,
and the product of these over $i$ is at least $e^{-q_{\max}^2T^2/N}$. If
$m_i\le n_i$ for every $i$, then all factors $1-j/n_i$ lie in $[0,1]$, and
their product is at least
$1-\sum_im_i^2/(2n_i)\ge1-M^2/(2n_{\min})$, where $n_{\min}=\min_in_i$. If
some $m_i>n_i$, then $\kappa_m=0$, while $M\ge n_{\min}+1$ gives
$M^2/(2n_{\min})\ge2$. In all cases
\begin{equation}\label{eq:A-kappa}
\Bigl(1-\frac{M^2}{2n_{\min}}\Bigr)e^{-q_{\max}^2T^2/N}\le\kappa_m\le e^{hq_{\max}M}.
\end{equation}

Let $w_m=v_m/\sum_{m'}v_{m'}$. This is a probability on $\Z^F_{\ge1}$,
since $H(t)>0$ under \textup{(H)}. By Lemma~\ref{lem:runs} the ratio in the
lemma is $\E_w\kappa_m$. So by~\eqref{eq:A-kappa} we need
$\E_we^{hq_{\max}M}$ and $\E_wM^2$. Since $W$ for $e^cA_F$ is
$e^{c(M-1)}W(m)$, we have, in the notation of the proof of
Lemma~\ref{lem:upper},
\[
\psi(c):=\log\E_we^{c(M-1)}=\log H_{e^c}(t)-\log H(t).
\]
The function $\psi$ is convex with $\psi(0)=0$, so $\psi(c)\le c\,\psi(1)$
for $0\le c\le1$. Also $\psi(1)\ge1$, since $M\ge k\ge2$. We show that
$\psi(1)\le C_{\mathcal K}T$. The bound for $H_\sigma$ in the proof of
Lemma~\ref{lem:upper} holds for every $\sigma>0$, and with $\sigma=e$ it
gives
$H_e(T\alpha)\le e^{k-1}b(\theta')\prod_{i\ge2}\theta_i\,e^{eT\langle\alpha,\theta\rangle}$
with $\theta=\theta(\alpha)$. For the lower bound, Theorem~\ref{thm:local}
gives $T_1$ such that, for $T\ge T_1$ and $\alpha\in\mathcal K$,
\[
H(T\alpha)=T^{-(k-1)}e^{T\langle q,\alpha\rangle}f_T(\alpha)
\ge\tfrac12(2\pi T)^{-\frac{k-1}2}(\mu\cdot r_\alpha)(l_\alpha\cdot\ones)
(\det\Sigma_\alpha)^{-1/2}e^{T\langle\alpha,\theta\rangle}.
\]
For $1\le T\le T_1$, the function $H(T\alpha)e^{-T\langle\alpha,\theta(\alpha)\rangle}$
is continuous and positive on the compact set $[1,T_1]\times\mathcal K$, so
it is bounded below there. The prefactor in the display is also continuous
and positive on $\mathcal K$. So in both cases
$H(T\alpha)\ge c_{\mathcal K}T^{-(k-1)/2}e^{T\langle\alpha,\theta\rangle}$.
Since $\theta(\alpha)$ and $b(\theta'(\alpha))$ are bounded on $\mathcal K$,
this gives $\psi(1)\le C_{\mathcal K}T$ for $T\ge1$.

Now let $x=hq_{\max}\le\frac12$. Then
$\E_we^{xM}=e^xe^{\psi(x)}\le e^{x(1+C_{\mathcal K}T)}=1+O(T^2/N)$, because
$x(1+C_{\mathcal K}T)=O(T^2/N)$ and $T^2/N\le1$. The Chernoff bound
$\Prob_w(M-1\ge y)\le e^{\psi(1)-y}$ gives
\[
\E_w(M-1)^2=\int_0^\infty2y\,\Prob_w(M-1\ge y)\,dy
\le\psi(1)^2+\int_{\psi(1)}^\infty2ye^{\psi(1)-y}\,dy
=\psi(1)^2+2\psi(1)+2,
\]
so $\E_wM^2=O(T^2)$. Finally $n_{\min}\ge\varepsilon_{\mathcal K}N$ with
$\varepsilon_{\mathcal K}=\min_{\alpha\in\mathcal K}\min_i\alpha_i>0$.
By~\eqref{eq:A-kappa},
\[
1-O\Bigl(\frac{T^2}N\Bigr)\le e^{-q_{\max}^2T^2/N}
\Bigl(1-\frac{\E_wM^2}{2\varepsilon_{\mathcal K}N}\Bigr)
\le\E_w\kappa_m\le\E_we^{hq_{\max}M}\le1+O\Bigl(\frac{T^2}N\Bigr),
\]
uniformly for $\alpha\in\mathcal K$ and $1\le T\le N^{1/2}$. \qed
